\documentclass[12pt,a4paper]{article}
\usepackage[margin=25mm, headheight=15pt, headsep=10pt]{geometry}
\usepackage{fontspec}
\usepackage[table]{xcolor} % Note: table option is required for rowcolors
\usepackage{titlesec}
\usepackage{tocloft}
\usepackage{fancyhdr}
\usepackage{booktabs}
\usepackage{tcolorbox}
\tcbuselibrary{skins, breakable}
\usepackage[colorlinks=true, linkcolor=emerald, urlcolor=emerald, bookmarks=true]{hyperref}

\definecolor{emerald}{RGB}{0,102,51}
\definecolor{darkred}{RGB}{139,0,0}
\definecolor{lightgray}{RGB}{245,245,245}

\usepackage[nil]{babel}
\babelprovide[import, main]{bengali}
\babelprovide[import]{arabic}
\babelfont{rm}[Renderer=HarfBuzz,Scale=1.0]{Noto Serif Bengali}
\babelfont{sf}[Renderer=HarfBuzz]{Noto Sans Bengali}
\babelfont{tt}[Renderer=HarfBuzz]{Noto Sans Bengali}
\babelfont[arabic]{rm}[Renderer=HarfBuzz,Scale=1.1]{Noto Naskh Arabic}

% Section styling
\titleformat{\section}{\Large\bfseries\color{emerald}}{\thesection.}{0.5em}{}
\titleformat{\subsection}{\large\bfseries\color{emerald!85!black}}{\thesubsection.}{0.5em}{}
\titleformat{\subsubsection}{\normalsize\bfseries\color{emerald!70!black}}{\thesubsubsection.}{0.5em}{}
\titlespacing*{\section}{0pt}{18pt}{8pt}

% Fancy Header/Footer setup
\pagestyle{fancy}
\fancyhf{} % clear all header and footer fields
\fancyhead[L]{\color{emerald}\small\textbf{\nouppercase{\leftmark}}}
\fancyfoot[L]{\scriptsize\color{black!50}{\lat{AI}}-সংকলিত, আলেম-যাচাইকৃত নয়}
\fancyfoot[C]{\color{emerald!70!black}\thepage}
\renewcommand{\headrulewidth}{0.4pt}
\renewcommand{\headrule}{\hbox to\headwidth{\color{emerald}\leaders\hrule height \headrulewidth\hfill}}
\renewcommand{\footrulewidth}{0pt}

% Custom boxes
% 1. Ayat/Hadith Box
\newtcolorbox{ayatbox}[1][]{
    enhanced, breakable, 
    colback=emerald!5, 
    colframe=emerald, 
    boxrule=0.5pt, 
    arc=3pt, 
    left=10pt, right=10pt, top=10pt, bottom=10pt,
    #1
}

% 2. Quote/Translation Box
\newtcolorbox{quotebox}[1][]{
    enhanced, breakable,
    frame hidden,
    colback=lightgray,
    borderline west={3pt}{0pt}{emerald},
    left=15pt, right=10pt, top=10pt, bottom=10pt,
    sharp corners,
    #1
}

% 3. AI-generated notice box
\newtcolorbox{warnbox}[1][]{
    enhanced, breakable,
    colback=red!4,
    colframe=darkred,
    boxrule=0.8pt,
    arc=3pt,
    left=10pt, right=10pt, top=10pt, bottom=10pt,
    #1
}

% Commands
\newcommand{\arab}[1]{\foreignlanguage{arabic}{#1}}
\newcommand{\kib}[1]{\textcolor{emerald}{\textbf{#1}}}

% Latin (English) fallback: Noto Serif Bengali has NO Latin glyphs
\newfontfamily\latfont[Renderer=HarfBuzz]{Noto Serif}
\newcommand{\lat}[1]{{\latfont #1}}

\setlength{\parskip}{5pt}
\setlength{\parindent}{0pt}

% Fix for missing bullet in Noto Serif Bengali
\renewcommand{\labelitemi}{$\bullet$}



\begin{document}
\begin{center}{\Large\bfseries\color{emerald} আল্লাহর রহমত ও তাওবা — ব্যবহারিক গাইড}\end{center}
\vspace{1em}
\section*{ব্যবহারিক গাইড — হতাশা নয়, ফিরে আসা: তাওবার পথ ও যাত্রা}

\begin{warnbox}
 \textbf{গুরুত্বপূর্ণ নোটিশ (\lat{Important} \lat{Notice}):} এই কন্টেন্টটি \lat{AI} (কৃত্রিম বুদ্ধিমত্তা) দিয়ে প্রস্তুত ও সংকলিত। \lat{AI} কঠোর সূত্র-মান নীতি মেনে শুধু নির্ভরযোগ্য উৎস থেকে তথ্য সংগ্রহ করেছে — কেবল সহীহ/হাসান হাদিস ও সহীহ সনদের আসার রাখা হয়েছে, দুর্বল সনদ বাদ দেওয়া হয়েছে। তবে এটি কোনো আলেম বা মুহাদ্দিস কর্তৃক যাচাইকৃত নয়।
\end{warnbox}

\begin{quote}
এই ফাইলটি তত্ত্বের পাশাপাশি \textbf{\lat{practice} (ব্যবহারিক)} উত্তর দেয় — \lat{“}আমি অনেক পাপ করেছি, এখন কী করব?\lat{”} — প্রতিটি উত্তরের পাশে কুরআন-হাদিসের দলিল। তাওবার শর্তাবলি উলামাদের প্রতিষ্ঠিত (\lat{established}) বক্তব্য থেকে নেওয়া — নতুন কোনো নিয়ম তৈরি করা হয়নি।
\end{quote}

\medskip\hrule\medskip

\subsection*{১. মূল বার্তা এক লাইনে (\lat{The} \lat{Core} \lat{Message})}

\begin{quote}
\textbf{পাপী মানুষের কোনো \lat{“}অতীত\lat{”} নেই আল্লাহর কাছে — আছে শুধু \lat{“}বর্তমানের ফেরা\lat{”}। তাওবার দরজা খোলা, গ্রহণকারী আল্লাহ নিজে, আর ফিরে এলে পুরনো পাপের বদলে লেখা হয় নতুন জীবন (ফুরকান ২৫:৭০)।}
\end{quote}

যে-কোনো গভীর প্রশ্নের জবাবে এই বার্তাটি ফিরে আসবে।

\medskip\hrule\medskip

\subsection*{২. কেন আল্লাহ বারবার তাওবার দরজা খোলা রেখেছেন?}

\begin{itemize}
\item \textbf{তাওবা গ্রহণ আল্লাহর নিজের গুণ:} তিনি \lat{“}তাওওয়াব\lat{”} — বারবার গ্রহণকারী (তাওবা ৯:১০৪)। নামটিই বলে — একবারের জন্য নয়, বারবারের জন্য দরজা।
\item \textbf{পাপ করা মানুষের স্বভাব, তাওবা তার প্রতিষেধক:} হাদিস — \lat{“}প্রতিটি আদম সন্তান গুনাহ করে, সেরা গুনাহকারী তাওবাকারী\lat{”} (ইবনে মাজাহ ৪২৫১)। আল্লাহ জানেন বান্দা দুর্বল; তাই তাওবার ব্যবস্থা রেখেছেন।
\item \textbf{তাওবা না হলে বান্দা হতাশায় ডুবে যেত:} শয়তানের প্রধান অস্ত্র হতাশা; দরজা খোলা রাখাই হতাশার প্রতিষেধক (যুমার ৩৯:৫৩)।
\item \textbf{আল্লাহ তাওবায় খুশি হন:} মরুভূমিতে হারানো উট ফিরে পাওয়ার আনন্দের উপমা (বুখারি ৬৩০৮) — এমন কোনো রব, যিনি বান্দার ফেরায় খুশি হন, তিনি দরজা বন্ধ রাখবেন কেন?
\end{itemize}
\textbf{উত্তর:} দরজা খোলা রাখা আল্লাহর নিজের ঘোষণা — এটা বান্দার \lat{“}সৌভাগ্য\lat{”} নয়, আল্লাহর গুণের ফল।

\medskip\hrule\medskip

\subsection*{৩. \lat{“}আমি অনেক পাপ করেছি — আল্লাহ আর ক্ষমা করবেন না\lat{”} — এই চিন্তার ইসলামি জবাব}

এই চিন্তাটি \textbf{তিনটি কুরআনি সূত্রে ভুল প্রমাণিত}:

\begin{enumerate}
\item \textbf{যুমার ৩৯:৫৩:} \lat{“}বলো: হে আমার বান্দারা, যারা নিজেদের উপর সীমালঙ্ঘন করেছ — তোমরা আল্লাহর রহমত থেকে নিরাশ হয়ো না। নিশ্চয়ই আল্লাহ \textbf{সমস্ত} গুনাহ ক্ষমা করে দেন।\lat{”} — সাহাবিরা একে কুরআনের সবচেয়ে আশাপূর্ণ আয়াত বলেছেন (ইবনে কাসীর)। \lat{“}সমস্ত\lat{”} — কোনো পাপ বাদ নেই (শিরকে মৃত্যু ছাড়া, নিসা ৪:৪৮)।
\item \textbf{তিরমিযী ৩৫৪০ (হাদিস কুদসি):} \lat{“}তোমার গুনাহ আকাশের মেঘ পর্যন্ত পৌঁছালেও… আমি ক্ষমা করব। পৃথিবী ভরতি পাপ নিয়ে এলে — পৃথিবী ভরতি ক্ষমা নিয়ে আসব।\lat{”}
\item \textbf{মুসলিম ২৭৫৯:} আল্লাহ রাতে-দিনে হাত বাড়ান — \lat{“}পশ্চিমে সূর্য উদিত হওয়া পর্যন্ত\lat{”}।
\end{enumerate}
\textbf{অতএব:} \lat{“}ক্ষমা করা হবে না\lat{”} — এটা আল্লাহ সম্পর্কে খারাপ ধারণা (সু'য যান্ন) — যা কুরআন নিষেধ করেছে (ফুসসিলাত ৪১:২৩-এর ভাবনায়) এবং পথভ্রষ্টদের বৈশিষ্ট্য বলেছে (হিজর ১৫:৫৬)।

\begin{quote}
 \textbf{ভারসাম্য:} এর মানে পাপে নির্ভীক হওয়া নয়। একই কুরআন সতর্ক করে — ৪:৯৩ (মুমিন হত্যা), ৪:৪৮ (শিরকে মৃত্যু)। আশা রাখবে — কিন্তু আশা মানে \lat{“}পাপ চালিয়ে যাওয়া\lat{”} নয়; আশা মানে \lat{“}ফিরে আসার সাহস\lat{”}।
\end{quote}

\medskip\hrule\medskip

\subsection*{৪. খাঁটি তাওবা (তাওবাতুন নাসূহ) আসলে কী?}

উলামারা তাওবার শর্তগুলো তাফসীর ও হাদিস থেকে স্থির করেছেন। সারমর্ম (\lat{established} \lat{conditions}):

\begin{table}[h]\centering\small
\begin{tabular}{lll}
শর্ত & বিবরণ & দলিলের ভিত্তি \\
\hline
\textbf{১. পাপ ছেড়ে দেওয়া} & সাথে সাথে পাপ বন্ধ — তাওবা মানে পাপ চালিয়ে যাওয়া নয় & ৩:১৩৫ — \lat{“}জেনে-বুঝে পাপে অটল থাকে না\lat{”} \\
\textbf{২. অনুতাপ (নাদম)} & পাপ করার জন্য অন্তরের আফসোস & \lat{“}অনুতাপই তাওবা\lat{”} (ইবনে মাজাহ ৪২৫২) \\
\textbf{৩. ফিরে না আসার সংকল্প} & ভবিষ্যতে পাপ না করার দৃঢ় নিয়ত & ৬৬:৮ — নসূহ তাওবার ব্যাখ্যায় উমার (রাঃ)-এর বক্তব্য (তাফসীর) \\
\textbf{৪. আল্লাহর হক আদায়} & ছুটে যাওয়া ফরজ/ওয়াজিবের কাযা ও ক্ষমা প্রার্থনা & তাফসীর ইবনে কাসীর — আলী (রাঃ)-এর শর্তের সারমর্ম \\
\textbf{৫. মানুষের হক ফেরত} & কারো ক্ষতি করলে — ক্ষমা/ক্ষতিপূরণ/ঋণ পরিশোধ & মুসলিম ২৫৭৭ — জুলুম নিষিদ্ধ; হাদিসে \lat{“}জুলুমের হক\lat{”} \\
\textbf{৬. সময়ের মধ্যে} & মৃত্যু কণ্ঠাগত হওয়ার আগে (তিরমিযী ৩৫৩৭) ও পশ্চিমে সূর্যোদয়ের আগে (মুসলিম ২৭৫৯) & — \\
\hline
\end{tabular}
\end{table}

\begin{quote}
\textbf{মূল কথা:} তাওবা = পাপ বন্ধ + আফসোস + সংকল্প + (হক থাকলে) আদায়। এটুকুই — কোনো জটিল আচার নেই। নবীজি (সা.) বলেছেন — তাওবা সহজ; যেখানে দাঁড়িয়ে আছ, সেখানেই আল্লাহকে ভয় করো, পাপের পর ভালো কাজ করো (তিরমিযী ১৯৮৭)।
\end{quote}

\medskip\hrule\medskip

\subsection*{৫. অতীতের পাপ কি ভবিষ্যৎ নির্ধারণ করে? (\lat{Does} \lat{the} \lat{past} \lat{decide} \lat{the} \lat{future}?)}

\textbf{না।} এর কোনো দলিল নেই; বরং বিপরীতের দলিল আছে:

\begin{itemize}
\item \textbf{ফুরকান ২৫:৭০:} তাওবাকারীর \textbf{পাপকে আল্লাহ নেকিতে বদলে দেন} — অতীতের হিসাব বদলে যায়; কিয়ামতে তাওবাকারী তার পাপের স্থানে নেকি দেখবে (ইবনে আব্বাস, ইবনে মাসউদের ব্যাখ্যা — তাফসীর)।
\item \textbf{৯৯ জন হত্যাকারী (মুসলিম ২৭৬৬):} তার অতীত ছিল ১০০ হত্যা — তারপরও আল্লাহর রহমত; ভবিষ্যৎ নির্ধারণ করল তার \textbf{শেষ অন্তরের ফেরা}, অতীত নয়।
\item \textbf{বেশ্যা ও কুকুর (বুখারি ৩৩২১):} অতীতের যিনা — একটি ইখলাসের ভালো কাজে সব মুছে গেল।
\item \textbf{কাব ইবনে মালিক (বুখারি ৪৪১৮):} ৫০ দিনের বিচ্ছিন্নতা ও পাপ-বোধের পর — আল্লাহ তাকে এমন মর্যাদায় বসালেন যে তার তাওবার ঘটনা কুরআনে (৯:১১৮) চিরকাল রইল। পাপীর অতীত আল্লাহর কাছে \lat{“}ফিরে আসার\lat{”} গল্পে রূপান্তরিত হলো।
\end{itemize}
\textbf{উত্তর:} ভবিষ্যৎ নির্ধারণ করে \textbf{আজকের ফেরা} — অতীতের পাপ নয়। শয়তানই চায় তুমি ভাবো \lat{“}অতীতের কারণে আর লাভ নেই\lat{”} — যেন তুমি আল্লাহর রহমতের হিসাব নিজে কষে বসে থাকো।

\medskip\hrule\medskip

\subsection*{৬. শয়তান কীভাবে \lat{guilt}-কে (অপরাধবোধ) \lat{despair}-এ (হতাশা) পরিণত করে?}

পাপের \lat{guilt} (অপরাধবোধ) নিজে খারাপ নয় — এটি আল্লাহর দিকে ফেরানোর চাবি হতে পারে (আল-ইমরান ৩:১৩৫ — পাপের পর আল্লাহকে স্মরণ ও ক্ষমা চাওয়া)। শয়তানের কাজ — এই \lat{guilt}-কে ঘুরিয়ে হতাশায় পরিণত করা:

\textbf{শয়তানের কৌশল ও কুরআন-হাদিসের জবাব:}

\begin{table}[h]\centering\small
\begin{tabular}{lll}
শয়তানের কৌশল & তার লক্ষ্য & জবাব (দলিল) \\
\hline
\lat{“}তোমার পাপ এত বড় — আল্লাহ ক্ষমা করবেন না\lat{”} & তাওবা থেকে বিরত রাখা & যুমার ৩৯:৫৩ — সমস্ত গুনাহ ক্ষমা; তিরমিযী ৩৫৪০ — মেঘ পর্যন্ত পাপ \\
\lat{“}এত দেরি হয়ে গেছে — আর লাভ নেই\lat{”} & ফেরার সাহস কেড়ে নেওয়া & তিরমিযী ৩৫৩৭ — প্রাণ কণ্ঠাগত হওয়ার আগ পর্যন্ত দরজা খোলা \\
\lat{“}আবার পাপ করলে তাওবা ভাঙবে — তাই এখনই করো না, পরে করব\lat{”} & ফেরা স্থগিত করা & মুসলিম ২৭৫৯ — বারবার পাপ-তাওবা কবুল; স্থগিত করা মানে মৃত্যুকে ডাকা \\
\lat{“}তুমি তো একদিনও ভালো নও — ধর্ম তোমার জন্য নয়\lat{”} & আত্মপরিচয় ভাঙা & মুসলিম ২৭৪৯ — পাপীই ক্ষমা চায়; মুসলিম ২৫৭৭ — \lat{“}তোমরা পাপ করো, আমি ক্ষমা করি\lat{”} \\
\lat{“}লোকেরা তোমাকে নিয়ে কী ভাববে — তুমি তো ধার্মিক ছিলে না\lat{”} & লজ্জায় ঢেকে রাখা & তাওবা ৯:১০২ — সাহাবিরাও পাপ স্বীকার করেছেন; আল্লাহর কাছে ফেরার লজ্জা নেই \\
\lat{“}আল্লাহ তোমার উপর রাগ করছেন — দোয়াও কবুল হবে না\lat{”} & দোয়া বন্ধ করা & বাকারা ২:১৮৬ — \lat{“}আমি নিকটে\lat{”}; গাফির ৪০:৬০ — \lat{“}ডাকো, আমি সাড়া দেব\lat{”} \\
\hline
\end{tabular}
\end{table}

\textbf{মূল চাবি:} শয়তানের সব কৌশলের লক্ষ্য একটাই — \textbf{তাওবার আগের মুহূর্তটিকে আটকে রাখা}। হাদিসে এসেছে — শয়তান প্রতিশ্রুতি দেয় আর ভুলিয়ে দেয় (নিসা ৪:১২০); আর আল্লাহর প্রতিশ্রুতি সত্য। যতক্ষণ শ্বাস আছে, শয়তানের কথা শোনার সময় শেষ হয়নি — আল্লাহর ডাক শোনার সময় আছে।

\medskip\hrule\medskip

\subsection*{৭. পাপের পর কীভাবে আবার শুরু করবেন? (\lat{How} \lat{to} \lat{restart} \lat{after} \lat{sin})}

এটি কোনো মনগড়া ধাপ নয় — কুরআন-সুন্নাহয় থাকা পথের সারমর্ম:

\begin{enumerate}
\item \textbf{থামুন (\lat{Stop}):} পাপটি এখনই বন্ধ করুন — পরের \lat{“}একবার\lat{”} নেই। ৩:১৩৫ — পাপে অটল না থাকা।
\item \textbf{আফসোস করুন (\lat{Regret}):} অন্তরে পাপের জন্য বেদনা অনুভব করুন — এটাই তাওবার মূল (ইবনে মাজাহ ৪২৫২)।
\item \textbf{ক্ষমা চান (\lat{Seek} \lat{forgiveness}):} সরাসরি আল্লাহর কাছে — \lat{“}রাব্বিগফিরলি\lat{”} / ইস্তিগফার; লজ্জায় কাউকে বলা জরুরি নয় (আল্লাহর কাছে সরাসরি — নিসা ৪:১১০)। যিনা-হত্যার মতো অপরাধে ইসলামি আদালত ছাড়া স্বীকার করার প্রয়োজন নেই — বরং গোপন রাখা উত্তম।
\item \textbf{ভালো কাজ দিয়ে পাপ ঢেকে দিন (\lat{Erase} \lat{with} \lat{good}):} নামাজ, দান, নফল — \lat{“}ভালো কাজ মন্দ কাজ মুছে দেয়\lat{”} (হুদ ১১:১১৪; তিরমিযী ১৯৮৭; তিরমিযী ২৬১৬)।
\item \textbf{পরিবেশ বদলান (\lat{Change} \lat{the} \lat{environment}):} পাপের জায়গা/সঙ্গী/সময় এড়িয়ে চলুন — ১০০ হত্যাকারীকে আলেম বলেছিলেন \lat{“}মন্দ জায়গা ছেড়ে ভালো জায়গায় যাও\lat{”} (মুসলিম ২৭৬৬)। পাপের ট্রিগার (ট্রিগার) দূর করা তাওবার অংশ।
\item \textbf{মানুষের হক আদায় করুন (\lat{Return} \lat{rights}):} কারো ক্ষতি/ঋণ থাকলে — ক্ষমা চান বা আদায় করুন (মুসলিম ২৫৭৭)।
\item \textbf{আবার পাপ হলে আবার তাওবা করুন (\lat{Repeat}):} পাপে ফিরে গেলে হতাশ হবেন না — আবার তাওবা; দরজা খোলা (মুসলিম ২৭৫৯)।
\end{enumerate}
\begin{quote}
\textbf{সবচেয়ে গুরুত্বপূর্ণ:} এক নম্বর ও দুই নম্বর — থামা ও আফসোস — এ দুটি হলেই তাওবা শুরু; বাকিগুলো তাওবার ফল ও স্থায়িত্বের জন্য।
\end{quote}

\medskip\hrule\medskip

\subsection*{৮. আল্লাহ কীভাবে ফিরে-আসা বান্দাকে সাহায্য করেন?}

কুরআন-হাদিসে আল্লাহর সাহায্যের কয়েকটি রূপ:

\begin{enumerate}
\item \textbf{হেদায়েত:} \lat{“}যারা আমার পথে সংগ্রাম করে, আমি অবশ্যই তাদের আমার পথ দেখাই\lat{”} (আনকাবুত ২৯:৬৯) — ফেরার চেষ্টাই আল্লাহর পথ দেখানোর শর্ত।
\item \textbf{নৈকট্য:} \lat{“}আমার দিকে এক বিঘত এগোল — আমি এক হাত এগিয়ে আসি…\lat{”} (মুসলিম ২৬৭৫) — বান্দার ছোট পদক্ষেপে আল্লাহ বড় সাড়া দেন।
\item \textbf{তাওবার সুযোগ:} \lat{“}আল্লাহ হাত বাড়ান…\lat{”} (মুসলিম ২৭৫৯) — এমনকি ফেরার \lat{“}ইচ্ছা\lat{”} ও \lat{“}সুযোগ\lat{”}ও আল্লাহর দান।
\item \textbf{ভালো সঙ্গী:} যুবকের হৃদয় পরিশুদ্ধ — নবীজির (সা.) দোয়ায় (মুসনাদ আহমাদ ২২২১৩); ভালো মানুষের দোয়া পাপীর পথ খুলে দেয়।
\item \textbf{পাপকে নেকিতে বদল:} ফুরকান ২৫:৭০ — তাওবার পরের আমলগুলোর মর্যাদা; পুরনো পাপের হিসাব মুছে নতুন নেকি লেখা।
\item \textbf{দোয়ার জবাব:} ইউনুসের দোয়া — \lat{“}কোনো মুসলিম তা দিয়ে দোয়া করলে আল্লাহ সাড়া দেন\lat{”} (তিরমিযী ৩৫০৫)।
\end{enumerate}
\textbf{ব্যবহারিক দিক:} ফিরে আসার চেষ্টা মানে — পাপ বন্ধের চেষ্টা, নামাজে ফেরা, ইস্তিগফারের অভ্যাস, ভালো মানুষের মজলিস, দোয়া। এই চেষ্টাগুলোর প্রতিটিই আল্লাহর সাহায্যের দরজা (আনকাবুত ২৯:৬৯)।

\medskip\hrule\medskip

\subsection*{৯. পাপের ভয় ও রহমতের আশা — ভারসাম্য (\lat{Balance} \lat{of} \lat{Fear} \lat{and} \lat{Hope})}

\begin{table}[h]\centering\small
\begin{tabular}{lll}
দিক & দলিল & ব্যবহার \\
\hline
\textbf{আশা (রজা)} & যুমার ৩৯:৫৩; তিরমিযী ৩৫৪০; মুসলিম ২৭৫৯ & হতাশা থেকে বাঁচতে — ফেরার সাহস জোগাতে \\
\textbf{ভয় (খওফ)} & নিসা ৪:৪৮, ৪:৯৩; গাফির ৪০:৩ (\lat{“}কঠোর শাস্তিদাতা\lat{”}) & পাপে নির্ভীক হওয়া থেকে বাঁচতে — সাবধান হতে \\
\hline
\end{tabular}
\end{table}

\textbf{উলামাদের পথ (\lat{established}):} রোগীর অবস্থা অনুযায়ী ওষুধ — পাপে ডুবে থাকা ব্যক্তির বেশি দরকার আশা; আত্মতুষ্ট/নির্ভীক ব্যক্তির বেশি দরকার ভয়। শেষ মুহূর্ত পর্যন্ত আশা ছাড়া যায় না, কিন্তু আশা পাপের অজুহাত হতে পারে না (তাফসীর ইবনে কাসীর — ৪০:৩-এর ব্যাখ্যা)।

\textbf{একটি সতর্ক দৃষ্টান্ত:} ৪:৯৩-এ ইচ্ছাকৃত মুমিন হত্যার ভয়াবহ সতর্কবাণী; ইবনে আব্বাস (রাঃ) বলেছেন — এর ব্যাখ্যায় \lat{“}তার তাওবা কবুল নয়\lat{”} (বুখারি ৪৭৬২)। অন্যদিকে জমহুর (সংখ্যাগরিষ্ঠ) — ২৫:৭০ ও মুসলিম ২৭৬৬-এর ভিত্তিতে — বলেন তাওবা কবুল হতে পারে, তবে নিহতের হক (কিসাস) থেকে যায়, পরিবার ক্ষমা করলে কিয়ামতে হিসাব হালকা হয় (ইবনে কাসীরের মীমাংসা)। অর্থ — আল্লাহর হক তাওবায় মাফ হয়, মানুষের হক মানুষের ক্ষমায়। এই মতভেদটি সৎভাবে জানা জরুরি — কারণ \lat{“}হত্যা করেছি, তাওবা করব\lat{”} ভাবতে গিয়ে কেউ যেন হত্যাকে হালকা না নেয়।

\medskip\hrule\medskip

\subsection*{১০. তাওবার পর আল্লাহর সাথে সম্পর্ক কীভাবে \lat{rebuild} (পুনর্গঠন) হয়?}

\begin{enumerate}
\item \textbf{নামাজ:} পাঁচ ওয়াক্ত নামাজ — পাপ মোছার মাধ্যম (হুদ ১১:১১৪; বুখারি ৫২৮-এর ভাবনা)। তাওবাকারীর প্রথম কাজ — নামাজে ফেরা।
\item \textbf{ইস্তিগফার ও জিকির:} সকাল-সন্ধ্যার জিকির, ইস্তিগফারের অভ্যাস — সম্পর্কের \lat{daily} (দৈনিক) সংযোগ।
\item \textbf{তাওবার নামাজ (সালাতুত তাওবা):} হাদিসে এসেছে — পাপ করলে ওযু করে দুই রাকাত নামাজ পড়া ও ক্ষমা চাওয়া (আবু দাউদ ১৫২১ — হাসান; তিরমিযী ৪০৬ — হাসান)।
\item \textbf{ভালো আমলের ধারা:} দান, সিলাতুর রাহিম, মানুষের উপকার — সাদাকা পাপ নিভায় (তিরমিযী ২৬১৬); এতে অন্তর আল্লাহর দিকে ঝোঁকে।
\item \textbf{দোয়া:} ইউনুসের দোয়া (তিরমিযী ৩৫০৫), ক্ষমার দোয়া — \lat{“}আল্লাহুম্মা ইন্নাকা আফুওউন তুহিব্বুল আফওয়া ফা'ফু আন্নি\lat{”} (তিরমিযী ৩৫১৩ — সহীহ)।
\item \textbf{সঙ্গ-সঙ্গী:} ভালো মানুষদের সাথে — পরিবেশ বদল তাওবার স্থায়িত্বের চাবি (মুসলিম ২৭৬৬)।
\end{enumerate}
\textbf{সম্পর্কের ধারা:} পাপের কারণে আল্লাহর সাথে দূরত্ব হয়েছে — তাওবা সেই দূরত্ব ভাঙে; তারপর নিয়মিত ইবাদত সম্পর্ককে টিকিয়ে রাখে। একদিনের তাওবা নয় — তাওবা একটি \lat{journey} (যাত্রা), যার প্রতিদিনের রুটিন হলো ইবাদত ও ইস্তিগফার।

\medskip\hrule\medskip

\subsection*{১১. যেসব কথা মনে রাখবেন (\lat{Reminder} \lat{Cards})}

\begin{quote}
\textbf{১. হতাশা হারাম} — \lat{“}আল্লাহ আমাকে ক্ষমা করবেন না\lat{”} ভাবা শিরক-পরবর্তী সবচেয়ে বড় ভুল ধারণা (যুমার ৩৯:৫৩)। \\
\textbf{২. তাওবা সহজ} — পাপ বন্ধ + আফসোস + সংকল্প; কোনো জটিল আচার নেই (ইবনে মাজাহ ৪২৫২)। \\
\textbf{৩. দরজা এখনো খোলা} — যতক্ষণ শ্বাস আছে, দরজা আছে (তিরমিযী ৩৫৩৭)। \\
\textbf{৪. আবার পাপ হলে — আবার তাওবা} — দরজা বারবার খোলা (মুসলিম ২৭৫৯)। \\
\textbf{৫. অতীত আল্লাহর হিসাবে নেই} — পাপের স্থানে নেকি লেখা হয় (ফুরকান ২৫:৭০)। \\
\textbf{৬. মানুষের হক আলাদা} — তাওবা আল্লাহর হক মাফ করায়; মানুষের ক্ষতি করলে তার ক্ষমা/ক্ষতিপূরণ দরকার (মুসলিম ২৫৭৭)। \\
\textbf{৭. আশা ও আমল দুই-ই} — শুধু মুখে তাওবা নয়; ভালো কাজ পাপ মুছে দেয় (তিরমিযী ১৯৮৭)। \\
\textbf{৮. শয়তান চায় থামাতে} — তার প্রতিটি কথা হতাশার দিকে নিয়ে যায়; আল্লাহর প্রতিটি ডাক ফেরার দিকে (নিসা ৪:১২০)।
\end{quote}

\medskip\hrule\medskip

\subsection*{১২. দোয়ার তালিকা (\lat{Du'a})}

\begin{table}[h]\centering\small
\begin{tabular}{lll}
দোয়া & অর্থ & উৎস \\
\hline
\textbf{রাব্বিগফিরলি ওয়া তুব আলাইয়্যা} & হে আমার রব! আমাকে ক্ষমা করো এবং আমার তাওবা কবুল করো & আবু দাউদ ১৫১২ — হাসান \\
\textbf{আল্লাহুম্মা ইন্নাকা আফুওউন তুহিব্বুল আফওয়া ফা'ফু আন্নি} & হে আল্লাহ! তুমি ক্ষমাশীল, ক্ষমা ভালোবাসো — আমাকে ক্ষমা করো & তিরমিযী ৩৫১৩ — সহীহ \\
\textbf{লা ইলাহা ইল্লা আনতা সুবহানাকা ইন্নি কুনতু মিনাজ জালিমিন} & তুমি ছাড়া কোনো ইলাহ নেই — তুমি পবিত্র; আমি জালিমদের অন্তর্ভুক্ত & তিরমিযী ৩৫০৫ — হাসান-সহীহ (ইউনুসের দোয়া) \\
\textbf{আস্তাগফিরুল্লাহা ওয়া আতূবু ইলাইহ} & আমি আল্লাহর কাছে ক্ষমা চাই এবং তাঁর দিকে তাওবা করি & মুসলিম ২৭০২-এর ভাবনা \\
\textbf{আল্লাহুম্মা ত্বাহহির ক্বলবী} & হে আল্লাহ! আমার অন্তর পবিত্র করো & মুসনাদ আহমাদ ২২২১৩-এর ভাবনা (নবীজির দোয়া) \\
\hline
\end{tabular}
\end{table}

\medskip\hrule\medskip

\subsection*{১৩. রেফারেন্স}

\begin{itemize}
\item কুরআন: যুমার ৩৯:৫৩, ফুরকান ২৫:৭০, তাওবা ৯:১০২, ৯:১০৪, ৯:১১৮, নিসা ৪:৪৮, ৪:৯৩, ৪:১১০, আল-ইমরান ৩:১৩৫-১৩৬, হুদ ১১:১১৪, আনকাবুত ২৯:৬৯, বাকারা ২:১৮৬, গাফির ৪০:৬০, ৬৬:৮, ইউসুফ ১২:৮৭, হিজর ১৫:৫৬
\item হাদিস: মুসলিম ২৫৭৭, ২৬৭৫, ২৭৪৯, ২৭৫৯, ২৭৬৬, ২৭৫৮; বুখারি ৬৩০৮, ৭৫০৭, ৩৩২১, ৪৭৬২; তিরমিযী ১৯৮৭, ২৬১৬, ৩৫৩৭, ৩৫৪০, ৩৫০৫, ৩৫১৩, ৪০৬; ইবনে মাজাহ ৪২৫১, ৪২৫২; আবু দাউদ ১৫২১, ১৫১২; মুসনাদ আহমাদ ২২২১৩
\item ব্যাখ্যা: তাফসীর ইবনে কাসীর (৩৯:৫৩, ২৫:৭০, ৬৬:৮, ৪:৯৩-এর ব্যাখ্যা), তাফসীর আস-সা'দী
\item এই ফাইলটি \lat{AI}-সংকলিত, আলেম-যাচাইকৃত নয় — গুরুত্বপূর্ণ সিদ্ধান্তের আগে একজন বিশ্বস্ত আলেমের সাথে \lat{consult} (পরামর্শ) করুন।
\end{itemize}

\end{document}
